% Language: LaTeX
% =====================================================================================================
%  paquetes.tex  --  SOLO carga de paquetes (sin configuracion)
%  La configuracion vive en formato.tex / listings.tex / entornos.tex
% =====================================================================================================

% --- Codificacion e idioma ---
\usepackage[utf8]{inputenc}
\usepackage[spanish]{babel}

% --- Tipografia ---
\usepackage{newtxtext,newtxmath}   % Times New Roman
\usepackage{ragged2e}              % justificado fino
\usepackage[normalem]{ulem}        % subrayado/tachado

% --- Color (antes de listings y de cualquier \definecolor) ---
\usepackage{xcolor}

% --- Geometria y espaciado (la CONFIG va en formato.tex) ---
\usepackage{geometry}
\usepackage{setspace}
\usepackage{lscape}                % paginas apaisadas

% --- Estructura, titulos, indices ---
\usepackage{titlesec}
\usepackage{enumitem}
\usepackage{tocbibind}
\usepackage{caption}
\usepackage{subcaption}
\usepackage{fancyhdr}

% --- Tablas ---
\usepackage{array}
\usepackage{booktabs}
\usepackage{tabularx}
\usepackage{longtable}
\usepackage{pbox}

% --- Graficos y codigo fuente ---
\usepackage{graphicx}
\usepackage{adjustbox}
\usepackage{verbatim}
\usepackage{listings}

% --- Matematicas ---
\usepackage{amsmath}
\let\Bbbk\relax                    % evita choque amssymb <-> newtxmath
\usepackage{amssymb}
\usepackage{bytefield}

% --- Cajas (tcolorbox) ---
% [most] no carga 'listings' con seguridad -> la pedimos explicita.
\usepackage[most]{tcolorbox}
\tcbuselibrary{listings,skins,breakable}

% --- Posicionamiento flotantes ---
\usepackage{float}                 % entorno [H]
\usepackage{newfloat}              % entornos flotantes propios (float "codigo")

% --- Enlaces: hyperref va casi al final ---
\usepackage[colorlinks=true, linkcolor=black, citecolor=blue]{hyperref}
\usepackage{url}
\usepackage{doi}                   % despues de hyperref

% --- Acronimos
\usepackage{acro}

% --- Etiquetas personalizadas TODO's
\usepackage{todonotes}
% comandos personalizados con colores
\newcommand{\TODO}[1]{\todo[color=green!40,inline]{TODO: #1}}
\newcommand{\REMOVE}[1]{\todo[color=red!40,inline]{REMOVE: #1}}
\newcommand{\REVIEW}[1]{\todo[color=yellow!40,inline]{REVIEW: #1}}



% \usepackage{hypcap}